\subsection{A cubic descent and the exact remaining cubic defect}

\begin{statement}{D23.F22}{Proved}{The index-three half-character is a descended cubic frame sign}
\end{statement}
There is a canonical continuous endomorphism
\[
 \mathcal D_E:\mathcal H\longrightarrow\mathcal H
\]
of the multiplier-one actual gentle group with the following explicit
description. In $B=\langle s,r\mid s^2=r^3=1\rangle^{\wedge}$ put
$\tau=sr$, $x=\tau^{-2}$, $y=sxs$, and
\[
 E=\langle r,srs\rangle\cong C_3\amalg C_3,
 \qquad C_f=f(x,r),\qquad t_f=\frac{3\Psi_3(f)}2.
\]
For $h=(1,f)\in\mathcal H$ first form the automorphism
\[
 \psi_f(r)=r,\qquad
 \psi_f(s)=f(r,srs)^{-1}s.
\]
Then
\[
 \boxed{\quad
 \mathcal D_E(h)=\operatorname{Inn}(x^{-t_f}C_f^{-1})\psi_f,
 \qquad A_{\mathcal D_E(h)}=x^{-t_f}C_f^{-1}.
 \quad}                                                     \tag{D23.F22.1}
\]
In particular
\[
 \boxed{\quad
 a_3(\mathcal D_E(h))=\Psi_3(h)/2\pmod2,\qquad
 \rho(\mathcal D_E(h))=\rho(h)\pmod3.
 \quad}                                                     \tag{D23.F22.2}
\]
Thus the cubic comparison defect of \textup{A.103}, on the actual
full-pentagon subgroup, has the intrinsic form
\[
 e_3=a_3+a_3\circ\mathcal D_E.
                                                               \tag{D23.F22.3}
\]
The fixed-point condition has the exact primitive formulation
\[
 \mathcal D_E(h)=h
 \quad\Longleftrightarrow\quad
 A_h^{-1}=f(x,r)x^{3\Psi_3(h)/2}.                              \tag{D23.F22.4}
\]
This is the equianharmonic primitive form with a specifically
normalized coefficient. Neither the fixed-point condition nor the
weaker equality $a_3=a_3\circ\mathcal D_E$ is asserted on all full
profinite GT here.

\begin{proof}
The kernel of the map $B\to C_2$ sending $s$ to the nonidentity
element and $r$ to one is $E=C_3\amalg C_3$, with the displayed
generators. This follows directly by the two-coset Reidemeister
rewriting, and the presentation holds profinitely as well. The
standard pure subgroup $F$ has index three in $E$; $r$ splits
$E\to C_3$, and conjugation by $r$ cyclically permutes
$x,y,(xy)^{-1}$.

The input automorphism $X\mapsto X$, $Y\mapsto f(X,Y)^{-1}Yf(X,Y)$
of the free profinite group preserves the closed normal subgroup
generated by $X^3,Y^3$, as does its inverse. Hence it induces an
automorphism of $E$, with $X=r,Y=srs$. Put $g=f(r,srs)$.
The two-cycle identity gives $sgs=g^{-1}$, so $(g^{-1}s)^2=1$.
Consequently the displayed formulas define an endomorphism $\psi_f$
of $B$. Its restriction to $E$ is exactly the induced automorphism,
and it induces the identity on $B/E$. It is therefore an
automorphism: surjectivity and injectivity follow successively from
the quotient and the restriction. The pair-composition rule
$f_{hk}=f_h\Phi_h(f_k)$ proves
$\psi_{f_{hk}}=\psi_{f_h}\psi_{f_k}$ by evaluating it in $E$.
In particular $h\mapsto\psi_f$ is a continuous homomorphism.

Both $g$ and $C_f$ belong to $F$: in $E/F=C_3$ the image of any
commutator word is one. Hence $\psi_f$ induces the identity on
$B/F=S_3$. The peripheral identity of \textup{A.100}, applied to
the pair $(r,srs)$, gives
\[
 \psi_f(y)=f(y,r)y f(y,r)^{-1},
 \qquad y=(rs)^{-2}.
\]
As $r^{-1}yr=x$ and $\psi_f(r)=r$, this gives
$\psi_f(x)=C_fxC_f^{-1}$.
It also gives $\psi_f(\tau)=C_f\tau C_f^{-1}$, with no choice
of square root. Indeed $v=C_f^{-1}\psi_f(\tau)C_f$ satisfies
$v^2=\tau^2$ and has the same image as $\tau$ in $S_3$.
The element $v\tau^{-1}\in F$ centralizes $x=\tau^{-2}$.
The pure cusp centralizer is $C_F(x)=\overline{\langle x\rangle}$;
thus $v=\tau x^b$ for some $b\in\widehat{\mathbb Z}$.
The equality of squares gives $2b=0$, whence $b=0$.

We compute the required integral normalization rather than choosing
it arbitrarily. Identify the $r$-module $F^{\rm ab}$ with
$\widehat{\mathbb Z}[\omega]$, where $[x]=1$, $[y]=\omega$,
$\omega^2+\omega+1=0$, and $r$ acts by multiplication by $\omega$.
In $F^{\rm ab}\rtimes\langle r\rangle$ the elements $r,x,srs$
are respectively
\[
 (0,\omega),\qquad(1,1),\qquad(-1,\omega^{-1}).
\]
The argument of \textup{A.90} works over the entire coefficient
ring $\widehat{\mathbb Z}[\omega]$: $1-\omega$ is a non-zero-divisor
because its norm is $3$, which is a non-zero-divisor in
$\widehat{\mathbb Z}$. The peripheral identity therefore gives
$A_f(\omega,\omega^{-1})=0$. The two-cycle and three-cycle then
give that the translation of
$f((0,\omega),(1,1))$ is a scalar $t\in\widehat{\mathbb Z}$.
Explicitly the two nontrivial translations in the three-cycle are
$-\bar t$ and $t$, so $t=\bar t$; the fixed ring is
$\widehat{\mathbb Z}$, again because multiplication by
$\omega-\omega^{-1}$ has nonzero norm $3$.

If $m_0,m_1,m_2$ are the three index-three Schreier coefficients,
the translation is $t=m_0+m_1\omega+m_2\omega^2$, with
$m_0+m_1+m_2=0$. Taking its trace gives
$2t=3m_0=3\Psi_3$. Since \textup{A.90} proves that $\Psi_3$ is
even, $t=t_f=3\Psi_3/2$.
The two-cycle now gives
\[
 C_f^{\rm ab}=-t_f[x],\qquad g^{\rm ab}=0.                    \tag{D23.F22.5}
\]
For the second assertion, evaluate the already vanishing expression
at $(0,\omega),(-1,\omega^{-1})$; changing the translation by a
scalar factor changes a commutator-word translation by the same
factor.

Let $\theta_f=\operatorname{Inn}(C_f^{-1})\psi_f$. It fixes
$\tau$, and its normalized pair coordinate is
\[
 f_0=s\theta_f(s)=(sC_f^{-1}s)gC_f.
\]
Conjugation by $s$ exchanges $[x]$ and $[y]$, so (D23.F22.5)
gives $f_0^{\rm ab}=(-t_f,t_f)$.
Conjugating once more by $x^{-t_f}$ changes its pair coordinate to
\[
 f_E=y^{-t_f}f_0x^{t_f}\in F'.
\]
The resulting automorphism fixes $\tau$, induces the identity on
$S_3$, and has commutator pair coordinate; equivalently it is an
actual multiplier-one gentle automorphism. It sends $r$ to its
conjugate by $x^{-t_f}C_f^{-1}\in F$, so uniqueness of the pure
cubic frame proves (D23.F22.1).

There is only one such normalized gentle representative in this
pure inner class. If two differ by $\operatorname{Inn}(u)$ with
$u\in F$, then fixing $\tau$ forces $u\in C_F(x)=\langle x\rangle$.
Writing $u=x^n$ changes a commutator pair coordinate by the
abelian vector $(-n,n)$, so $n=0$. The homomorphism law for
$\psi_f$ and this uniqueness prove the homomorphism law for
$\mathcal D_E$. Its formula and the continuity of $\Psi_3$ give
continuity.

The sign computation must retain the signed lift. Put
\[
 U=\begin{pmatrix}1&1\\0&1\end{pmatrix},\quad
 L=\begin{pmatrix}1&0\\-1&1\end{pmatrix},\quad
 S=ULU,\quad R=-UL
 \qquad\hbox{over }\mathbb F_3.
\]
Define $q_E:E\to\operatorname{SL}_2(\mathbb F_3)$ by
$r\mapsto R$ and $srs\mapsto SRS^{-1}$. Both images have order
three, so this is well-defined. Direct multiplication gives
\[
 q_E(x)=-U^2,\qquad q_E(y)=-L^2.
\]
Consequently $q_E(w)=(-I)^{\ell(w)}q_3(w)$ on $F$, where
$\ell(w)$ is the sum modulo two of its two abelian coordinates,
and $q_3$ is the positive signed free-subgroup lift of \textup{AD.2}.
In particular $q_E=q_3$ on the pure cubic frame in (D23.F22.1),
whose abelianization is zero by (D23.F22.5).

Write $\rho=\rho(h)\bmod3$ and
$M_\rho=\left(\begin{smallmatrix}1&\rho\\\rho&1+\rho^2\end{smallmatrix}\right)$.
The supplied signed covariance calculation gives
$f(U^2,L^2)=M_\rho$. Since $f$ has zero exponent sums, a central
sign in either input does not change its value. Using
$R=U^2L^2U^{-2}$ gives
\[
 C_E:=q_E(C_f)=f(-U^2,R)=U^2M_\rho U^{-2}.                   \tag{D23.F22.6}
\]
Put $b=\Psi_3(h)/2$, so $t_f=3b$. Because $U^3=I$,
\[
 q_3(A_{\mathcal D_E(h)})=q_E(A_{\mathcal D_E(h)})
       =(-I)^bC_E^{-1}.
\]
For $\rho_D=\rho(\mathcal D_E(h))\bmod3$, the corrected
cubic frame is therefore
\[
 T_{\mathcal D_E(h)}=(-I)^b C_E U^{2\rho_D}.
\]
It centralizes $R$ by \textup{AD.2}. Exact multiplication for
$\rho=0,1,2$ gives
$C_EU^{2\rho}=R^\rho$. Since no nonidentity power of $U$
centralizes $R$, this forces $\rho_D=\rho$, and hence
\[
 T_{\mathcal D_E(h)}=(-I)^bR^\rho.
\]
The sign character of $\langle-I\rangle\times\langle R\rangle$
is precisely $a_3$, proving (D23.F22.2). All signed identities,
all three shear residues, and both values of $b\bmod2$ are
independently checkable in \texttt{cubic\_descent.py}; its record
is \texttt{cubic\_descent.json}.

On full GT, \textup{A.76} identifies $a_3$ with
$\kappa_d\bmod2$, and \textup{A.103} defines
$e_3=\kappa_d+\Psi_3/2$. This proves (D23.F22.3).
Finally two normalized gentle automorphisms fixing $\tau$ and
having the same pure $r$-frame coincide: their images of $r$
coincide and $s=\tau r^{-1}$. Applying (D23.F22.1) proves
(D23.F22.4).
\end{proof}

\paragraph{Source boundary.}
The normalizations, actual gentle group, pure-frame uniqueness, and
the positive signed covariance are the supplied \textup{R3--R4},
\textup{LM7\_PUSH E.65}, and \textup{DELTA\_ON\_GT AD.1--AD.2}.
The character definition, evenness and peripheral identity are
\textup{A.88--A.90} and \textup{A.100}.
The new assertion here is the cubic endomorphism and its exact
parity formula. \textup{A.65}, equation \textup{(M8)}, records the
equianharmonic primitive form with an existential even coefficient;
the computation above identifies the coefficient for this descent.
The primary geometric account of the two fillings is
Lochak--Nakamura--Schneps, \emph{Eigenloci of 5 point configurations},
\S6, equations (6.3)--(6.4), author PDF pp.30--31, and its
Theorem 2, equation (HF1), author PDF p.4. That theorem has Galois
scope. Enriquez, \emph{Quasi-reflection algebras and cyclotomic
associators}, arXiv:math/0408035v2, Theorem 8.1 and equation (50),
pp.30--32, supplies the full-GT cyclic distribution relation but
is not silently substituted for the missing fixed-point or parity
comparison in (D23.F22.4).
